%%=============================================================================
%% De inhoud van je projectvoorstel
%%=============================================================================

%---------- 1. Context en probleemstelling ------------------------------------
\section{Context en probleemstelling}%
\label{vst:context}

\textbf{Context}
\begin{itemize}
    \item Een digitaal platform dat artsen ondersteunt 
bij consultaties en administratieve taken. 
\end{itemize}

\textbf{Probleemstelling}
\begin{itemize}
    \item De inhoud en opmaak van verstuurde e-mails
zijn onvoldoende aanpasbaar. 
\end{itemize}
Gevolg: communicatie tussen arts en patiënt 
voldoen niet aan beide partijen hun specifieke behoeftes. 

%---------- 2. Doelstelling ----------------------------------------------------
\section{Doelstelling}%
\label{vst:doelstelling}
\begin{itemize}
    \item Een proof of concept (PoC) ontwikkelen van een WYSIWYG-e-maileditor
    \item Gebruikers kunnen templates* aanpassen, opslaan, vooraf bekijken en als testmail versturen 
\end{itemize}
templates*: afspraakbevestiging, afspraakherinnering en afspraakannulering
 
\newpage
%---------- 3. Technologieën --------------------------------------------------
\section{Technologieën}%
\label{vst:technologieen}
\textbf{Frontend}
\begin{itemize}
    \item Angular 22 (opbouw van de interactieve UI)
    \item TypeScript(voor typed en onderhoud code)
    \item Quill 2  (basis voor de visuele WYSIWYG-editor)
\end{itemize}

\medskip
\textbf{Mock-API en opslag}
\begin{itemize}
    \item Node.js (uitvoering van de servercode)
    \item TypeScript(getypeerde serverlogica)
    \item Express(opzetten van API-endpoints voor het laden en opslaan van templates)
\end{itemize}

\medskip
\textbf{HTML-verwerking en e-mailrendering}
\begin{itemize}
    \item Handlebars(templateverwerking en het behoud van de vaste layoutschil)
    \item \texttt{sanitize-html}(om onveilige HTML te verwijderen)
    \item Placeholdervervanging met escaping (om dynamische waarden gecontroleerd en veilig in te voegen)
    \item Juice (om CSS inline te plaatsen voor een betere ondersteuning door e-mailclients)
    \item \texttt{html-to-text} (om naast HTML ook een tekstuele e-mailversie te genereren)
\end{itemize}

\medskip
\textbf{Testomgeving}
\begin{itemize}
    \item Docker (reproduceerbare lokale testomgeving)
    \item Mailpit (om uitgaande testmails lokaal op te vangen en te controleren)
\end{itemize}

\medskip
\textbf{Geautomatiseerde tests}
\begin{itemize}
    \item Vitest (unit- en integratietests)
    \item Playwright (browsercontroles en screenshotvergelijkingen)
\end{itemize}

\medskip
\textbf{Versiebeheer}
\begin{itemize}
    \item Bitbucket (versiebeheer, samenwerking en het bijhouden van codewijzigingen)
\end{itemize}


%---------- 4. Afbakening ------------------------------------------------------
\section{Afbakening}%
\label{vst:afbakening}
\textbf{Out-Of-Scope (OOS)}
\begin{itemize}
    \item Integratie in productieplatform
    \item Plannen of versturen van productie e-mails
    \item Keycloak e-mails met echte patiëntgegevens en AI-functionaliteiten
\end{itemize}
De bestaande layoutschil met logo, merkkleuren en footer blijft behouden. 
Een mogelijk integratievoorstel kan als de kernfunctionaliteiten en belangrijkste tests correct en tijdig afgerond zijn.

%---------- 5. Verwachte meerwaarde -------------------------------------------
\section{Verwachte meerwaarde}%
\label{vst:meerwaarde}

\textbf{Inzichten in}
\begin{itemize}
    \item Technische haalbaarheid, beperkingen en veiligheidsaspecten van visuele e-mailbewerking. 
\end{itemize}

\medskip  
\textbf{Deliverables}
\begin{itemize}
    \item Een werkende demonstratie
    \item Een reproduceerbare testopstelling 
    \item Geautomatiseerde tests en testresultaten van verschillende e-mailclients
\end{itemize}


%---------- 6. Planning en werkbelasting --------------------------------------
\section{Planning en werkbelasting}%
\label{vst:planning}

\begin{center}
\small
\begin{tabular}{p{0.36\linewidth}p{0.47\linewidth}r}
\toprule
\textbf{Fase} & \textbf{Resultaat} & \textbf{Uren} \\
\midrule
Analyse & Scope, vereisten en testcriteria & 12 \\
Technologiekeuze en opzet & Technologiekeuzes en ontwikkelomgeving & 10 \\
Ontwerp & Architectuur, editorconcept en testmatrix & 13 \\
Realisatie & Editor, templates, preview en testmails & 40 \\
Testen en validatie & Geautomatiseerde tests en clientcontroles & 15 \\
Rapport, poster en demo & Eindrapport en demonstratie & 10 \\
\midrule
\textbf{Totaal} & & \textbf{100} \\
\bottomrule
\end{tabular}
\end{center}

%---------- 7. Risico's --------------------------------------------------------
\section{Risico's}%
\label{vst:risicos}

\textbf{Beperkte werkbelasting}
\begin{itemize}
    \item Niet alle uitbreidingen worden tijdig gerealiseerd
\end{itemize}

\textbf{Verschillen tussen e-mailclients}
\begin{itemize}
    \item Gegenereerde inhoud en opmaak wordt verschillend weergegeven (Outlook)
    \item CSS inline 
    \item Output tests in Gmail, Outlook en Apple
    \item Outlook Desktop wordt meegenomen voor zover de beschikbare testomgeving het toelaat
\end{itemize}

\textbf{Verschillen met de bestaande backend}
\begin{itemize}
    \item De mock-API kan afwijken van de bestaande Laravel-backend van Vesalius.ai
    \item Verschillen in architectuur en gegevensverwerking
\end{itemize}

